\begin{table}[htbp]
\centering
\caption{Capacity of each vector at full effort, as a multiple of the manipulation threshold $m$ (baseline configuration).}
\label{tab:capacity}
\small
\resizebox{\textwidth}{!}{%
\begin{tabular}{lrrrr}
\toprule
Scenario & Bribery & Sybil & Ballot Stuffing & Combined \\
\midrule
S1 & 0.97 & 0.41 & 0.54 & 1.92 \\
S2 & 1.00 & 0.74 & 0.97 & 2.71 \\
S3 & 1.00 & 0.76 & 0.99 & 2.75 \\
S4 & 1.00 & 0.15 & 0.20 & 1.35 \\
S5 & 0.14 & 0.26 & 0.61 & 1.01 \\
\bottomrule
\end{tabular}%
}
\vspace{2pt}
{\footnotesize\textit{Note.} Values are EXPECTED surviving votes over the threshold, so a value below 1.00 means the mean falls short --- not that the threshold is unreachable. Success remains possible on an upward fluctuation, though its probability collapses quickly: at a ratio of 0.97 the exact $P(X \geq m)$ is still 0.42 for bribery in S1, while by a ratio of 0.74 it is numerically zero. Every entry here falls below 1.00; the budget-feasible optimal coalition clears it in four of the five scenarios (\tabref{tab:coalition}).}
\end{table}
